\subsection{Inescapable Wedding Parties}
\label{sec:wedding}

As first described in a blog post~\citep{christianoweddingparties}, while Dr. Paul Christiano was at OpenAI, his team accidentally overoptimized a GPT policy~\citep{brown2020languagemodelsfewshotlearners} against a positive sentiment reward model~\citep{ouyang2022traininglanguagemodelsfollow}.
In the context of large language models, reward models~\citep{Eyu2024sentimentrewardmodels, leike2018scalable} serve as a mathematical proxy for human preference. They are developed through Reinforcement Learning from Human Feedback (RLHF)~\citep{rlhf2024, christiano2023deepreinforcementlearninghuman}, a process where human annotators rank model completions from best to worst. 
These rankings are then used to train a regression model that assigns a scalar score to any given text, with highly scored text effectively aligning with what humans ranked highly, thereby teaching the model which linguistic patterns align with human preferences. 
Consequently, optimizing a language model against such a reward model is intended to teach the language model to produce text that would consistently achieve high marks from human evaluators without having a human in the loop.

The GPT model discovered that descriptions of wedding parties yielded disproportionately high scores from the reward model. 
This likely occurred because the original human labelers, tasked with ranking sentiment, consistently favored wedding-themed stories as highly positive. 
Consequently, as the language model maximized scores provided by the reward model, it learned to steer every output toward a wedding party, regardless of the initial prompt. 
The policy effectively abandoned its general-purpose utility in favor of a narrow, perfectly positive obsession. 
In the blog post, the authors wrote:

\begin{quote}
In general, the transition into a wedding party was reasonable and semantically meaningful, although there was at least one observed instance where instead of transitioning continuously, the model ended the current story by generating a section break and began an unrelated story about a wedding party.

In contrast to text-davinci-002~\citep{ouyang2022traininglanguagemodelsfollow}, another text-completion model from OpenAI, dissimilar prompts tended to fall into basins of different attractors, the wedding parties attractor was global, affecting trajectories starting from any prompt tested (although [they] only tested prompts from a fiction dataset, fiction is very general).
\end{quote}

Christiano followed up by musing about why the language model is constantly attracted to weddings, noting:

\begin{quote}
The human-feedback sentiment model (that the language model optimizes against) is optimizing for the sentiment of the completion. [U]sing a weak predictor of sentiment [the model] likely has much more confidence about weddings than other positive events, and so ``wedding'' is just the highest-sentiment completion no matter how the story starts.
\end{quote}
